\newcommand{\spacerule}[0]{\vspace{15pt}\hrule\vspace{15pt}}
\newcommand{\dblrule}{\vspace{5pt}\hrule\vspace{2pt}\hrule\vspace{5pt}}
\newcommand{\roundbox}[8]{\begin{mbbox}{#1}{#2}{#3}{#4}{#5}{#6} \color{#7}#8 \end{mbbox}}
\newcommand{\scaledfrac}[2]{\frac{\raisebox{1pt}{\scaleobj{0.8}{#1}}}{\raisebox{-4pt}{\scaleobj{0.8}{#2}}}}
\newcommand{\ceiling}[1]{\left \lceil #1 \right \rceil}
\newcommand{\floor}[1]{\left \lfloor #1 \right \rfloor}
\newcommand{\vecprod}[1]{\Pi\parenths{\boldsymbol{#1}}}
\newcommand{\smily}[1]{\overset{\lower0.25em\hbox{$\smash{\scriptscriptstyle\smile}$}}{#1}}
\newcommand{\frowny}[1]{\overset{\lower0.25em\hbox{$\smash{\scriptscriptstyle\frown}$}}{#1}}
\newcommand{\plusy}[1]{\overset{\lower0.25em\hbox{$\smash{\scriptscriptstyle +}$}}{#1}}
\newcommand{\circly}[1]{\overset{\lower0.25em\hbox{$\smash{\scriptscriptstyle \circ}$}}{#1}}
\newcommand{\overy}[2]{\overset{\lower0.25em\hbox{$\smash{\scriptscriptstyle #2}$}}{#1}}
\newcommand{\highlight}[1]{\fontfamily{qag}\textcolor{red}{\textbf{\emph{#1}}}\fontfamily{lmr}}
\newcommand{\subhighlight}[1]{\fontfamily{qag}\textcolor{black}{\textbf{\emph{#1}}}\fontfamily{lmr}}
\newcommand{\parenths}[1]{\left ( #1 \right )}
\newcommand{\braces}[1]{\left \{ #1 \right \}}
\newcommand{\brackets}[1]{\left [ #1 \right ]}
\newcommand{\absolute}[1]{\left | #1 \right |}
\newcommand{\rightopen}[1]{\left [ #1 \right )}
\newcommand{\prob}[1]{\text{Pr} \braces{#1}}
\newcommand{\mean}[1]{\text{E} \brackets{#1}}
\newcommand{\meansq}[1]{\mean{\absolute{#1}^2}}
\newcommand{\mat}[1]{\brackets{\begin{matrix}#1\end{matrix}}}
\newcommand{\realop}[1]{\mathfrak{Re}\braces{#1}}
\newcommand{\imop}[1]{\mathfrak{Im}\braces{#1}}
\newcommand{\range}[3]{#1,#2,\cdots,#3}
\newcommand{\mbitem}[1]{
\begin{minipage}[t]{0.9\linewidth}
#1 \end{minipage}}
\newcommand{\tableheading}[1]{
\begin{minipage}[t]{1\linewidth}
\vspace{0.1cm}\center\textbf{#1}\vspace{0.1cm} \end{minipage}}
\newcommand{\tableitem}[1]{
\begin{minipage}[t]{1\linewidth}
\vspace{0.1cm}\color{black} #1\vspace{0.1cm} \end{minipage}}
\newcommand{\notascalar}{Italic symbols like $a$ and $M$ denote scalar quantities and functions}
\newcommand{\notavector}{Bold italic symbols like $\boldsymbol{a}$ and $\boldsymbol{M}$ denote vector quantities and functions}
\newcommand{\notamatrix}{Bold non-italic upper case symbols like $\textbf{M}$ denote matrices}
\newcommand{\notanonnegintegerset}{$\circly{\mathbb{I}}$ denotes the set of non-negative integer numbers}
\newcommand{\notaintegerset}{$\mathbb{I}$ denotes the set of integer numbers}
\newcommand{\notaevenintegerset}{$\mathbb{E}$ denotes the set of even integer numbers}
\newcommand{\notaoddintegerset}{$\mathbb{O}$ denotes the set of odd integer numbers}
\newcommand{\notaposintegerset}{$\plusy{\mathbb{I}}$ denotes the set of positive integer numbers}
\newcommand{\notaposevenintegerset}{$\plusy{\mathbb{E}}$ denotes the set of positive even integer numbers}
\newcommand{\notanonnegevenintegerset}{$\circly{\mathbb{E}}$ denotes the set of non-negative even integer numbers}
\newcommand{\notaposoddintegerset}{$\plusy{\mathbb{O}}$ denotes the set of positive odd integer numbers}
\newcommand{\notaintegerpowerstwo}{$\mathbbm{2}^{\mathbb{I}}$ denotes the set of integer powers of two}
\newcommand{\notaposintegerpowerstwo}{$\mathbbm{2}^{\plusy{\mathbb{I}}}$ denotes the set of positive integer powers of two}
\newcommand{\notaposevenintegerpowerstwo}{$\mathbbm{2}^{\plusy{\mathbb{E}}}$ denotes the set of positive even integer powers of two}
\newcommand{\notanonnegevenintegerpowerstwo}{$\mathbbm{2}^{\circly{\mathbb{E}}}$ denotes the set of non-negative even integer powers of two}
\newcommand{\notaposoddintegerpowerstwo}{$\mathbbm{2}^{\plusy{\mathbb{O}}}$ denotes the set of positive odd integer powers of two}
\newcommand{\notanonnegintegerpowerstwo}{$\mathbbm{2}^{\circly{\mathbb{I}}}$ denotes the set of non-negative integer powers of two}
\newcommand{\notaposintegerpowersfour}{$\mathbbm{4}^{\plusy{\mathbb{I}}}$ denotes the set of positive integer powers of four}
\newcommand{\notaintegersetkl}{$\mathbb{I}_{k,l}$, where $k\le l$ and $k,l\in\mathbb{I}$, denotes the set of integer numbers from $k$ to $l$}
\newcommand{\notaintegersetN}{$\mathbb{I}_N$, where $N\in\plusy{\mathbb{I}}$, denotes the set of integer numbers from $0$ to $N-1$. Note that $\mathbb{I}_N=\mathbb{I}_{0,N-1}$}
\newcommand{\notaposrealset}{$\plusy{\mathbb{R}}$ denotes the set of positive real numbers}
\newcommand{\notanonnegrealset}{$\circly{\mathbb{R}}$ denotes the set of non-negative real numbers}
\newcommand{\notarealset}{$\mathbb{R}$ denotes the set of real numbers}
\newcommand{\notaopenrealinterval}{Open real interval $\parenths{a,b}$ denotes the set of all values $x\in\mathbb{R}$ such that $a<x<b$}
\newcommand{\notaclosedrealinterval}{Closed real interval $\brackets{a,b}$ denotes the set of all values $x\in\mathbb{R}$ such that $a\le x\le b$}
\newcommand{\notarightopenrealinterval}{Right-open real interval $\rightopen{a,b}$ denotes the set of all values $x\in\mathbb{R}$ such that $a\le x< b$}
\newcommand{\notaconjugate}{Superscript $^*$ denotes complex conjugation}
\newcommand{\notatranspose}{Superscript $^T$ denotes matrix transposition}
\newcommand{\notacomplexset}{$\mathbb{C}$ denotes the set of complex numbers}
\newcommand{\notafloor}{$\floor{\cdot}$ denotes the floor of (i.e., the largest integer that is smaller than or equal to) the enclosed real quantity}
\newcommand{\notamean}{$\mean{\cdot}$ denotes mathematical expectation (mean value) of the enclosed random variable}
\newcommand{\notaprob}{$\prob{\cdot}$ denotes probability of the enclosed event}
\newcommand{\notarealpart}{$\realop{\cdot}$ denotes the real part of the enclosed quantity}
\newcommand{\notaimaginarypart}{$\imop{\cdot}$ denotes the imaginary part of the enclosed quantity}

\definecolor{titles}{HTML}{CC00FF}
\definecolor{subtitles}{HTML}{E97132}
\definecolor{math}{HTML}{0066FF}
\definecolor{DarkRosyBrown}{RGB}{176, 84, 48}
\definecolor{BGPink}{RGB}{244, 224, 244}
\definecolor{BGBrown}{RGB}{228, 204, 184}
\definecolor{BGGreen}{HTML}{CAEEF3}
\definecolor{BGBlue}{HTML}{66FFFF}
\definecolor{BrightGreen}{RGB}{44, 164, 48}

\newenvironment{colorequ}[1]{\equation\color{#1}}{\endequation}
\newenvironment{coloritemize}{\itemize\color{blue}}{\enditemize}
\newenvironment{colorsubitemize}{\itemize\color{magenta}}{\enditemize}
\newenvironment{colorsubsubitemize}{\itemize\color{teal}}{\enditemize}
\newcommand{\coloritem}{\item[\ding{88}]}
\newcommand{\colorsubitem}{\item[\ding{70}]}
\newcommand{\colorsubsubitem}{\item[\ding{228}]}

\newcommand{\indentedsubsub}[2]{
\vspace{10pt}\subsubsection{#1}\hfill\\

\vspace{-15pt}#2
}

\newcounter{DraftBox}
\newcommand{\DraftBox}[2]{\refstepcounter{DraftBox}\begin{tcolorbox}  [colback=orange!5!white,colframe=orange!75!black,title=\fontfamily{cmr}\scshape\textbf{#1.\theDraftBox}] #2\end{tcolorbox}}

\newcounter{QuestionBox}
\newcommand{\QuestionBox}[1]{\refstepcounter{QuestionBox}\begin{tcolorbox}  [colback=white,colframe=blue,title=\fontfamily{cmr}Q.\theQuestionBox- #1] \end{tcolorbox}}

\usepackage{tabularx,makecell,booktabs,multirow}

\usepackage{adjustbox} % if not already loaded
\usepackage{pgfplots}
\pgfplotsset{compat=1.17}



\setcounter{dbltopnumber}{3}
\renewcommand{\dbltopfraction}{.99}
\setcounter{topnumber}{3}
\renewcommand{\topfraction}{.99}
\setcounter{totalnumber}{5}
\renewcommand{\textfraction}{.01}
\setlength{\textfloatsep}{6pt plus 1pt minus 2pt}
\makeatletter\setlength{\@fptop}{0pt}\makeatother

\usepackage{pifont}    % For \xmark

\newcommand{\cmark}{\ding{51}}%
\newcommand{\xmark}{\ding{55}}%

\usepackage{amssymb}   % for \checkmark
\usepackage{makecell}  % for multi-line cells with \makecell